\documentclass{article}%
\usepackage[T1]{fontenc}%
\usepackage[utf8]{inputenc}%
\usepackage{lmodern}%
\usepackage{textcomp}%
\usepackage{lastpage}%
\usepackage{geometry}%
\geometry{margin=1in,headheight=14pt,headsep=25pt}%
\usepackage{hyperref}%
\usepackage{enumitem}%
\usepackage{fancyhdr}%
\usepackage{xcolor}%
\usepackage{xurl}%
\usepackage{breakurl}%
\usepackage{amsmath}%
\usepackage{amssymb}%
\usepackage{mathtools}%
\usepackage{amsthm}%
\usepackage{thmtools}%
%

            \hypersetup{
                colorlinks=true,
                linkcolor=blue,
                filecolor=magenta,
                urlcolor=blue,
                breaklinks=true,
                bookmarks=true,
                pdfborder={0 0 0}
            }
            \urlstyle{same}
            \Urlmuskip=0mu plus 1mu  % URL breaking in nicer spots
            
            % Math configuration
            \everymath{\displaystyle}
            \setlength{\jot}{10pt}  % Increase spacing between equations
            
            \pagestyle{fancy}
            \fancyhf{}
            \rhead{Generated on \today}
            \lhead{Course Summary}
            \cfoot{\thepage}
            
            % Better Unicode support
            \usepackage[utf8]{inputenc}
            \usepackage{newunicodechar}
            
            % Configure itemize settings globally
            \setlist[itemize]{leftmargin=*,nosep}
        %
\title{Artifical Intelligence Basics Course Summary}%
\author{Generated by Scribe.AI}%
\date{\today}%
%
\begin{document}%
\normalsize%
\maketitle%
\section{Key Terms}%
\label{sec:KeyTerms}%
\subsection{Fundamental AI Concepts}%
\label{subsec:FundamentalAIConcepts}%
\begin{itemize}[leftmargin=*]%
\item \href{\url{https://purdue.brightspace.com/d2l/common/assets/pdfjs-d2l-dist/1.0.14-legacy/web/viewer.html?file=\%252Fcontent\%252Fenforced\%252F1095467-wl.202510.CS.24300.001\%252F02L2-Intro.pdf\%253Fou\%253D1095467\&lang=en-us\&container=d2l-fileviewer-rendered-pdf\&fullscreen=d2l-fileviewer-rendered-pdf-dialog\&height=667}#page=22}{\textbf{Minimax Search}}: A search algorithm for two-player zero-sum games that aims to find the best move for the current player, assuming the opponent also plays optimally.%
\item \href{\url{https://purdue.brightspace.com/d2l/common/assets/pdfjs-d2l-dist/1.0.14-legacy/web/viewer.html?file=\%252Fcontent\%252Fenforced\%252F1095467-wl.202510.CS.24300.001\%252F02L2-Intro.pdf\%253Fou\%253D1095467\&lang=en-us\&container=d2l-fileviewer-rendered-pdf\&fullscreen=d2l-fileviewer-rendered-pdf-dialog\&height=667}#page=23}{\textbf{Alpha-Beta Pruning}}: A technique to reduce the search space in Minimax search by avoiding the evaluation of branches that cannot affect the optimal decision.%
\item \href{\url{https://purdue.brightspace.com/d2l/common/assets/pdfjs-d2l-dist/1.0.14-legacy/web/viewer.html?file=\%252Fcontent\%252Fenforced\%252F1095467-wl.202510.CS.24300.001\%252F02L2-Intro.pdf\%253Fou\%253D1095467\&lang=en-us\&container=d2l-fileviewer-rendered-pdf\&fullscreen=d2l-fileviewer-rendered-pdf-dialog\&height=667}#page=6}{\textbf{Turing Test}}: An operational test for intelligent behavior, proposed by Alan Turing, where a human evaluator must distinguish between a human and a machine based on their textual responses.%
\item \href{\url{https://purdue.brightspace.com/d2l/common/assets/pdfjs-d2l-dist/1.0.14-legacy/web/viewer.html?file=\%252Fcontent\%252Fenforced\%252F1095467-wl.202510.CS.24300.001\%252F02L2-Intro.pdf\%253Fou\%253D1095467\&lang=en-us\&container=d2l-fileviewer-rendered-pdf\&fullscreen=d2l-fileviewer-rendered-pdf-dialog\&height=667}#page=11}{\textbf{Game Tree}}: A tree-like structure representing all possible sequences of moves in a game, with leaf nodes representing the outcomes.%
\item \href{\url{https://purdue.brightspace.com/d2l/common/assets/pdfjs-d2l-dist/1.0.14-legacy/web/viewer.html?file=\%252Fcontent\%252Fenforced\%252F1095467-wl.202510.CS.24300.001\%252F02L2-Intro.pdf\%253Fou\%253D1095467\&lang=en-us\&container=d2l-fileviewer-rendered-pdf\&fullscreen=d2l-fileviewer-rendered-pdf-dialog\&height=667}#page=46}{\textbf{Reinforcement Learning}}: A type of machine learning where an agent learns to make decisions by interacting with an environment and receiving rewards or penalties.%
\item \href{\url{https://purdue.brightspace.com/d2l/common/assets/pdfjs-d2l-dist/1.0.14-legacy/web/viewer.html?file=\%252Fcontent\%252Fenforced\%252F1095467-wl.202510.CS.24300.001\%252F05L2-ML\%20Basics.pdf\%253Fou\%253D1095467\&lang=en-us\&container=d2l-fileviewer-rendered-pdf\&fullscreen=d2l-fileviewer-rendered-pdf-dialog\&height=667}#page=3}{\textbf{Inductive Reasoning}}: The process of deriving general principles from specific observations.%
\item \href{\url{https://purdue.brightspace.com/d2l/common/assets/pdfjs-d2l-dist/1.0.14-legacy/web/viewer.html?file=\%252Fcontent\%252Fenforced\%252F1095467-wl.202510.CS.24300.001\%252F05L2-ML\%20Basics.pdf\%253Fou\%253D1095467\&lang=en-us\&container=d2l-fileviewer-rendered-pdf\&fullscreen=d2l-fileviewer-rendered-pdf-dialog\&height=667}#page=6}{\textbf{Linear Regression}}: A machine learning algorithm used to model the relationship between a dependent variable and one or more independent variables by fitting a linear equation to observed data.%
\item \href{\url{https://purdue.brightspace.com/d2l/common/assets/pdfjs-d2l-dist/1.0.14-legacy/web/viewer.html?file=\%252Fcontent\%252Fenforced\%252F1095467-wl.202510.CS.24300.001\%252F05L2-ML\%20Basics.pdf\%253Fou\%253D1095467\&lang=en-us\&container=d2l-fileviewer-rendered-pdf\&fullscreen=d2l-fileviewer-rendered-pdf-dialog\&height=667}#page=8}{\textbf{Model Evaluation}}: The process of assessing the performance of a machine learning model.%
\item \href{\url{https://purdue.brightspace.com/d2l/common/assets/pdfjs-d2l-dist/1.0.14-legacy/web/viewer.html?file=\%252Fcontent\%252Fenforced\%252F1095467-wl.202510.CS.24300.001\%252F05L2-ML\%20Basics.pdf\%253Fou\%253D1095467\&lang=en-us\&container=d2l-fileviewer-rendered-pdf\&fullscreen=d2l-fileviewer-rendered-pdf-dialog\&height=667}#page=14}{\textbf{Overfitting}}: A phenomenon in machine learning where a model learns the training data too well, resulting in poor generalization to unseen data.%
\item \href{\url{https://purdue.brightspace.com/d2l/common/assets/pdfjs-d2l-dist/1.0.14-legacy/web/viewer.html?file=\%252Fcontent\%252Fenforced\%252F1095467-wl.202510.CS.24300.001\%252F05L2-ML\%20Basics.pdf\%253Fou\%253D1095467\&lang=en-us\&container=d2l-fileviewer-rendered-pdf\&fullscreen=d2l-fileviewer-rendered-pdf-dialog\&height=667}#page=15}{\textbf{Empirical Risk Minimizer}}: A machine learning approach that aims to find the model that minimizes the average loss on the training data.%
\item \href{\url{https://purdue.brightspace.com/d2l/common/assets/pdfjs-d2l-dist/1.0.14-legacy/web/viewer.html?file=\%252Fcontent\%252Fenforced\%252F1095467-wl.202510.CS.24300.001\%252F05L2-ML\%20Basics.pdf\%253Fou\%253D1095467\&lang=en-us\&container=d2l-fileviewer-rendered-pdf\&fullscreen=d2l-fileviewer-rendered-pdf-dialog\&height=667}#page=20}{\textbf{Binary Classification}}: A machine learning task where the goal is to classify data into one of two categories.%
\item \href{\url{https://purdue.brightspace.com/d2l/common/assets/pdfjs-d2l-dist/1.0.14-legacy/web/viewer.html?file=\%252Fcontent\%252Fenforced\%252F1095467-wl.202510.CS.24300.001\%252F05L2-ML\%20Basics.pdf\%253Fou\%253D1095467\&lang=en-us\&container=d2l-fileviewer-rendered-pdf\&fullscreen=d2l-fileviewer-rendered-pdf-dialog\&height=667}#page=24}{\textbf{Decision Boundary}}: A line or surface that separates different classes in a classification problem.%
\item \href{\url{https://purdue.brightspace.com/d2l/common/assets/pdfjs-d2l-dist/1.0.14-legacy/web/viewer.html?file=\%252Fcontent\%252Fenforced\%252F1095467-wl.202510.CS.24300.001\%252F05L2-ML\%20Basics.pdf\%253Fou\%253D1095467\&lang=en-us\&container=d2l-fileviewer-rendered-pdf\&fullscreen=d2l-fileviewer-rendered-pdf-dialog\&height=667}#page=25}{\textbf{0-1 Loss}}: A loss function in classification that assigns a loss of 1 if the prediction is incorrect and 0 if it is correct.%
\item \href{\url{https://purdue.brightspace.com/d2l/common/assets/pdfjs-d2l-dist/1.0.14-legacy/web/viewer.html?file=\%252Fcontent\%252Fenforced\%252F1095467-wl.202510.CS.24300.001\%252F05L2-ML\%20Basics.pdf\%253Fou\%253D1095467\&lang=en-us\&container=d2l-fileviewer-rendered-pdf\&fullscreen=d2l-fileviewer-rendered-pdf-dialog\&height=667}#page=27}{\textbf{Perceptron Loss}}: A loss function used in the perceptron algorithm that penalizes misclassifications.%
\item \href{\url{https://purdue.brightspace.com/d2l/common/assets/pdfjs-d2l-dist/1.0.14-legacy/web/viewer.html?file=\%252Fcontent\%252Fenforced\%252F1095467-wl.202510.CS.24300.001\%252F05L2-ML\%20Basics.pdf\%253Fou\%253D1095467\&lang=en-us\&container=d2l-fileviewer-rendered-pdf\&fullscreen=d2l-fileviewer-rendered-pdf-dialog\&height=667}#page=27}{\textbf{Hinge Loss}}: A loss function used in support vector machines that penalizes predictions that are not sufficiently confident.%
\item \href{\url{https://purdue.brightspace.com/d2l/common/assets/pdfjs-d2l-dist/1.0.14-legacy/web/viewer.html?file=\%252Fcontent\%252Fenforced\%252F1095467-wl.202510.CS.24300.001\%252F05L2-ML\%20Basics.pdf\%253Fou\%253D1095467\&lang=en-us\&container=d2l-fileviewer-rendered-pdf\&fullscreen=d2l-fileviewer-rendered-pdf-dialog\&height=667}#page=27}{\textbf{Logistic Loss}}: A loss function used in logistic regression that measures the difference between the predicted probability and the true label.%
\item \href{\url{https://purdue.brightspace.com/d2l/common/assets/pdfjs-d2l-dist/1.0.14-legacy/web/viewer.html?file=\%252Fcontent\%252Fenforced\%252F1095467-wl.202510.CS.24300.001\%252F05L2-ML\%20Basics.pdf\%253Fou\%253D1095467\&lang=en-us\&container=d2l-fileviewer-rendered-pdf\&fullscreen=d2l-fileviewer-rendered-pdf-dialog\&height=667}#page=28}{\textbf{Gradient Descent}}: An iterative optimization algorithm used to find the minimum of a function by following the direction of the negative gradient.%
\item \href{\url{https://purdue.brightspace.com/d2l/common/assets/pdfjs-d2l-dist/1.0.14-legacy/web/viewer.html?file=\%252Fcontent\%252Fenforced\%252F1095467-wl.202510.CS.24300.001\%252F05L2-ML\%20Basics.pdf\%253Fou\%253D1095467\&lang=en-us\&container=d2l-fileviewer-rendered-pdf\&fullscreen=d2l-fileviewer-rendered-pdf-dialog\&height=667}#page=29}{\textbf{Stochastic Gradient Descent}}: A variant of gradient descent that updates the model parameters using a randomly selected subset of the training data.%
\end{itemize}

%
\subsection{Neural Networks}%
\label{subsec:NeuralNetworks}%
\begin{itemize}[leftmargin=*]%
\item \href{\url{https://purdue.brightspace.com/d2l/common/assets/pdfjs-d2l-dist/1.0.14-legacy/web/viewer.html?file=\%252Fcontent\%252Fenforced\%252F1095467-wl.202510.CS.24300.001\%252F02L2-Intro.pdf\%253Fou\%253D1095467\&lang=en-us\&container=d2l-fileviewer-rendered-pdf\&fullscreen=d2l-fileviewer-rendered-pdf-dialog\&height=667}#page=45}{\textbf{Perceptron}}: A fundamental unit in neural networks that performs a weighted sum of its inputs and applies an activation function to produce an output.%
\end{itemize}

%
\subsection{Linear Algebra for AI}%
\label{subsec:LinearAlgebraforAI}%
\begin{itemize}[leftmargin=*]%
\item \href{\url{https://purdue.brightspace.com/d2l/common/assets/pdfjs-d2l-dist/1.0.14-legacy/web/viewer.html?file=\%252Fcontent\%252Fenforced\%252F1095467-wl.202510.CS.24300.001\%252F04L1-Linear\%20Algebra\%20Primer.pdf\%253Fou\%253D1095467\&lang=en-us\&container=d2l-fileviewer-rendered-pdf\&fullscreen=d2l-fileviewer-rendered-pdf-dialog\&height=667}#page=4}{\textbf{Matrix}}: A rectangular array of numbers, symbols, or expressions arranged in rows and columns.%
\item \href{\url{https://purdue.brightspace.com/d2l/common/assets/pdfjs-d2l-dist/1.0.14-legacy/web/viewer.html?file=\%252Fcontent\%252Fenforced\%252F1095467-wl.202510.CS.24300.001\%252F04L1-Linear\%20Algebra\%20Primer.pdf\%253Fou\%253D1095467\&lang=en-us\&container=d2l-fileviewer-rendered-pdf\&fullscreen=d2l-fileviewer-rendered-pdf-dialog\&height=667}#page=5}{\textbf{Linear Transform}}: A function that maps vectors from one vector space to another, satisfying $T(u + v) = T(u) + T(v)$ and $T(cu) = cT(u)$ for all vectors $u$, $v$ and scalars $c$.%
\item \href{\url{https://purdue.brightspace.com/d2l/common/assets/pdfjs-d2l-dist/1.0.14-legacy/web/viewer.html?file=\%252Fcontent\%252Fenforced\%252F1095467-wl.202510.CS.24300.001\%252F04L1-Linear\%20Algebra\%20Primer.pdf\%253Fou\%253D1095467\&lang=en-us\&container=d2l-fileviewer-rendered-pdf\&fullscreen=d2l-fileviewer-rendered-pdf-dialog\&height=667}#page=14}{\textbf{Matrix Multiplication}}: An operation that combines two matrices to produce a third matrix, where the element at row $i$ and column $j$ of the resulting matrix is the dot product of row $i$ of the first matrix and column $j$ of the second matrix. The inner dimensions must match.%
\item \href{\url{https://purdue.brightspace.com/d2l/common/assets/pdfjs-d2l-dist/1.0.14-legacy/web/viewer.html?file=\%252Fcontent\%252Fenforced\%252F1095467-wl.202510.CS.24300.001\%252F04L1-Linear\%20Algebra\%20Primer.pdf\%253Fou\%253D1095467\&lang=en-us\&container=d2l-fileviewer-rendered-pdf\&fullscreen=d2l-fileviewer-rendered-pdf-dialog\&height=667}#page=10}{\textbf{Matrix Inverse}}: For a square matrix $A$, its inverse $A^{-1}$ is a matrix such that $AA^{-1} = A^{-1}A = I$, where $I$ is the identity matrix.%
\item \href{\url{https://purdue.brightspace.com/d2l/common/assets/pdfjs-d2l-dist/1.0.14-legacy/web/viewer.html?file=\%252Fcontent\%252Fenforced\%252F1095467-wl.202510.CS.24300.001\%252F04L1-Linear\%20Algebra\%20Primer.pdf\%253Fou\%253D1095467\&lang=en-us\&container=d2l-fileviewer-rendered-pdf\&fullscreen=d2l-fileviewer-rendered-pdf-dialog\&height=667}#page=16}{\textbf{Matrix Transposition}}: An operation that swaps the rows and columns of a matrix. The transpose of matrix $A$ is denoted as $A^T$.%
\item \href{\url{https://purdue.brightspace.com/d2l/common/assets/pdfjs-d2l-dist/1.0.14-legacy/web/viewer.html?file=\%252Fcontent\%252Fenforced\%252F1095467-wl.202510.CS.24300.001\%252F04L1-Linear\%20Algebra\%20Primer.pdf\%253Fou\%253D1095467\&lang=en-us\&container=d2l-fileviewer-rendered-pdf\&fullscreen=d2l-fileviewer-rendered-pdf-dialog\&height=667}#page=21}{\textbf{Derivative}}: A measure of how a function changes with respect to a change in its input.%
\item \href{\url{https://purdue.brightspace.com/d2l/common/assets/pdfjs-d2l-dist/1.0.14-legacy/web/viewer.html?file=\%252Fcontent\%252Fenforced\%252F1095467-wl.202510.CS.24300.001\%252F04L1-Linear\%20Algebra\%20Primer.pdf\%253Fou\%253D1095467\&lang=en-us\&container=d2l-fileviewer-rendered-pdf\&fullscreen=d2l-fileviewer-rendered-pdf-dialog\&height=667}#page=19}{\textbf{Partial Derivative}}: A derivative of a function of multiple variables with respect to one of those variables, holding the others constant.%
\item \href{\url{https://purdue.brightspace.com/d2l/common/assets/pdfjs-d2l-dist/1.0.14-legacy/web/viewer.html?file=\%252Fcontent\%252Fenforced\%252F1095467-wl.202510.CS.24300.001\%252F04L1-Linear\%20Algebra\%20Primer.pdf\%253Fou\%253D1095467\&lang=en-us\&container=d2l-fileviewer-rendered-pdf\&fullscreen=d2l-fileviewer-rendered-pdf-dialog\&height=667}#page=21}{\textbf{Sigmoid Function}}: A mathematical function that maps any input value to an output value between 0 and 1. Often used as an activation function in neural networks.%
\end{itemize}

%
\section{Problem Types}%
\label{sec:ProblemTypes}%
\subsection{Fundamental AI Concepts}%
\label{subsec:FundamentalAIConcepts}%
\begin{itemize}[leftmargin=*]%
\item \href{\url{https://purdue.brightspace.com/d2l/common/assets/pdfjs-d2l-dist/1.0.14-legacy/web/viewer.html?file=\%252Fcontent\%252Fenforced\%252F1095467-wl.202510.CS.24300.001\%252F02L2-Intro.pdf\%253Fou\%253D1095467\&lang=en-us\&container=d2l-fileviewer-rendered-pdf\&fullscreen=d2l-fileviewer-rendered-pdf-dialog\&height=667}#page=5}{\textbf{Understanding the Nature of Intelligence}}: This problem explores the definition and characteristics of intelligence, both human and artificial. It examines various perspectives on what constitutes intelligence and how it can be measured or simulated.%
\item \href{\url{https://purdue.brightspace.com/d2l/common/assets/pdfjs-d2l-dist/1.0.14-legacy/web/viewer.html?file=\%252Fcontent\%252Fenforced\%252F1095467-wl.202510.CS.24300.001\%252F02L2-Intro.pdf\%253Fou\%253D1095467\&lang=en-us\&container=d2l-fileviewer-rendered-pdf\&fullscreen=d2l-fileviewer-rendered-pdf-dialog\&height=667}#page=5}{\textbf{Defining Artificial Intelligence}}: This problem involves formulating a precise definition of artificial intelligence, considering different perspectives and historical contexts. It explores the challenges in defining AI and the implications of different definitions.%
\item \href{\url{https://purdue.brightspace.com/d2l/common/assets/pdfjs-d2l-dist/1.0.14-legacy/web/viewer.html?file=\%252Fcontent\%252Fenforced\%252F1095467-wl.202510.CS.24300.001\%252F02L2-Intro.pdf\%253Fou\%253D1095467\&lang=en-us\&container=d2l-fileviewer-rendered-pdf\&fullscreen=d2l-fileviewer-rendered-pdf-dialog\&height=667}#page=6}{\textbf{Passing the Turing Test}}: This problem investigates the implications of passing the Turing Test as a measure of artificial intelligence. It examines whether passing the test truly indicates intelligence or merely the ability to mimic human behavior.%
\item \href{\url{https://purdue.brightspace.com/d2l/common/assets/pdfjs-d2l-dist/1.0.14-legacy/web/viewer.html?file=\%252Fcontent\%252Fenforced\%252F1095467-wl.202510.CS.24300.001\%252F02L2-Intro.pdf\%253Fou\%253D1095467\&lang=en-us\&container=d2l-fileviewer-rendered-pdf\&fullscreen=d2l-fileviewer-rendered-pdf-dialog\&height=667}#page=31}{\textbf{Similarities and Differences Between AI Tasks}}: This problem explores the commonalities and differences between seemingly disparate AI tasks, such as game playing and speech synthesis, to identify underlying principles and transferable techniques.%
\end{itemize}

%
\subsection{AI Algorithms \& Optimization}%
\label{subsec:AIAlgorithmsOptimization}%
\begin{itemize}[leftmargin=*]%
\item \href{\url{https://purdue.brightspace.com/d2l/common/assets/pdfjs-d2l-dist/1.0.14-legacy/web/viewer.html?file=\%252Fcontent\%252Fenforced\%252F1095467-wl.202510.CS.24300.001\%252F02L2-Intro.pdf\%253Fou\%253D1095467\&lang=en-us\&container=d2l-fileviewer-rendered-pdf\&fullscreen=d2l-fileviewer-rendered-pdf-dialog\&height=667}#page=11}{\textbf{Search Space Complexity in Game Playing}}: This problem addresses the computational challenges of searching game trees, particularly in complex games like chess and Go. It explores the limitations of brute-force search and the need for more efficient algorithms.%
\item \href{\url{https://purdue.brightspace.com/d2l/common/assets/pdfjs-d2l-dist/1.0.14-legacy/web/viewer.html?file=\%252Fcontent\%252Fenforced\%252F1095467-wl.202510.CS.24300.001\%252F02L2-Intro.pdf\%253Fou\%253D1095467\&lang=en-us\&container=d2l-fileviewer-rendered-pdf\&fullscreen=d2l-fileviewer-rendered-pdf-dialog\&height=667}#page=12}{\textbf{Overcoming Limitations of Search Algorithms}}: This problem focuses on the limitations of search-based approaches to AI and the need for incorporating machine learning techniques to handle the complexity of real-world problems.%
\item \href{\url{https://purdue.brightspace.com/d2l/common/assets/pdfjs-d2l-dist/1.0.14-legacy/web/viewer.html?file=\%252Fcontent\%252Fenforced\%252F1095467-wl.202510.CS.24300.001\%252F02L2-Intro.pdf\%253Fou\%253D1095467\&lang=en-us\&container=d2l-fileviewer-rendered-pdf\&fullscreen=d2l-fileviewer-rendered-pdf-dialog\&height=667}#page=50}{\textbf{Understanding the AlphaGo System}}: This problem examines the architecture and functionality of the AlphaGo system, highlighting its combination of deep learning and search algorithms.%
\item \href{\url{https://purdue.brightspace.com/d2l/common/assets/pdfjs-d2l-dist/1.0.14-legacy/web/viewer.html?file=\%252Fcontent\%252Fenforced\%252F1095467-wl.202510.CS.24300.001\%252F05L2-ML\%20Basics.pdf\%253Fou\%253D1095467\&lang=en-us\&container=d2l-fileviewer-rendered-pdf\&fullscreen=d2l-fileviewer-rendered-pdf-dialog\&height=667}#page=28}{\textbf{Gradient Descent}}: An iterative optimization algorithm used to minimize the loss function by following the gradient direction.%
\item \href{\url{https://purdue.brightspace.com/d2l/common/assets/pdfjs-d2l-dist/1.0.14-legacy/web/viewer.html?file=\%252Fcontent\%252Fenforced\%252F1095467-wl.202510.CS.24300.001\%252F05L2-ML\%20Basics.pdf\%253Fou\%253D1095467\&lang=en-us\&container=d2l-fileviewer-rendered-pdf\&fullscreen=d2l-fileviewer-rendered-pdf-dialog\&height=667}#page=29}{\textbf{Stochastic Gradient Descent}}: A computationally efficient variant of gradient descent that randomly selects data points to update the model's parameters.%
\end{itemize}

%
\subsection{AI Applications \& Challenges}%
\label{subsec:AIApplicationsChallenges}%
\begin{itemize}[leftmargin=*]%
\item \href{\url{https://purdue.brightspace.com/d2l/common/assets/pdfjs-d2l-dist/1.0.14-legacy/web/viewer.html?file=\%252Fcontent\%252Fenforced\%252F1095467-wl.202510.CS.24300.001\%252F02L2-Intro.pdf\%253Fou\%253D1095467\&lang=en-us\&container=d2l-fileviewer-rendered-pdf\&fullscreen=d2l-fileviewer-rendered-pdf-dialog\&height=667}#page=15}{\textbf{Addressing Challenges in Computer Vision}}: This problem examines the historical context and ongoing challenges in computer vision, highlighting the gap between initial expectations and the persistent difficulties in achieving human-level performance.%
\item \href{\url{https://purdue.brightspace.com/d2l/common/assets/pdfjs-d2l-dist/1.0.14-legacy/web/viewer.html?file=\%252Fcontent\%252Fenforced\%252F1095467-wl.202510.CS.24300.001\%252F02L2-Intro.pdf\%253Fou\%253D1095467\&lang=en-us\&container=d2l-fileviewer-rendered-pdf\&fullscreen=d2l-fileviewer-rendered-pdf-dialog\&height=667}#page=20}{\textbf{Ethical Considerations in AI}}: This problem addresses the ethical implications of AI systems, particularly concerning bias and fairness. It examines the potential for AI to perpetuate or amplify existing societal biases.%
\item \href{\url{https://purdue.brightspace.com/d2l/common/assets/pdfjs-d2l-dist/1.0.14-legacy/web/viewer.html?file=\%252Fcontent\%252Fenforced\%252F1095467-wl.202510.CS.24300.001\%252F02L2-Intro.pdf\%253Fou\%253D1095467\&lang=en-us\&container=d2l-fileviewer-rendered-pdf\&fullscreen=d2l-fileviewer-rendered-pdf-dialog\&height=667}#page=25}{\textbf{Challenges in Designing Autonomous Robots}}: This problem investigates the complexities of designing autonomous robots, considering the need for robust perception, planning, and decision-making capabilities in dynamic environments.%
\end{itemize}

%
\subsection{Machine Learning Fundamentals}%
\label{subsec:MachineLearningFundamentals}%
\begin{itemize}[leftmargin=*]%
\item \href{\url{https://purdue.brightspace.com/d2l/common/assets/pdfjs-d2l-dist/1.0.14-legacy/web/viewer.html?file=\%252Fcontent\%252Fenforced\%252F1095467-wl.202510.CS.24300.001\%252F02L2-Intro.pdf\%253Fou\%253D1095467\&lang=en-us\&container=d2l-fileviewer-rendered-pdf\&fullscreen=d2l-fileviewer-rendered-pdf-dialog\&height=667}#page=19}{\textbf{Vulnerability of Neural Networks to Adversarial Attacks}}: This problem explores the susceptibility of neural networks to adversarial attacks, where small, carefully crafted perturbations can lead to misclassifications.%
\item \href{\url{https://purdue.brightspace.com/d2l/common/assets/pdfjs-d2l-dist/1.0.14-legacy/web/viewer.html?file=\%252Fcontent\%252Fenforced\%252F1095467-wl.202510.CS.24300.001\%252F04L1-Linear\%20Algebra\%20Primer.pdf\%253Fou\%253D1095467\&lang=en-us\&container=d2l-fileviewer-rendered-pdf\&fullscreen=d2l-fileviewer-rendered-pdf-dialog\&height=667}#page=17}{\textbf{Understanding Perceptrons}}: A perceptron, a fundamental building block of neural networks, is introduced, illustrating its structure and function.%
\item \href{\url{https://purdue.brightspace.com/d2l/common/assets/pdfjs-d2l-dist/1.0.14-legacy/web/viewer.html?file=\%252Fcontent\%252Fenforced\%252F1095467-wl.202510.CS.24300.001\%252F04L1-Linear\%20Algebra\%20Primer.pdf\%253Fou\%253D1095467\&lang=en-us\&container=d2l-fileviewer-rendered-pdf\&fullscreen=d2l-fileviewer-rendered-pdf-dialog\&height=667}#page=19}{\textbf{Calculus in AI}}: The slide highlights the importance of calculus, specifically derivatives, partial derivatives, and finding maxima/minima/gradients, in AI algorithms.%
\item \href{\url{https://purdue.brightspace.com/d2l/common/assets/pdfjs-d2l-dist/1.0.14-legacy/web/viewer.html?file=\%252Fcontent\%252Fenforced\%252F1095467-wl.202510.CS.24300.001\%252F04L1-Linear\%20Algebra\%20Primer.pdf\%253Fou\%253D1095467\&lang=en-us\&container=d2l-fileviewer-rendered-pdf\&fullscreen=d2l-fileviewer-rendered-pdf-dialog\&height=667}#page=21}{\textbf{Calculating the Derivative of the Sigmoid Function}}: The slide shows the step-by-step derivation of the derivative of the sigmoid function using calculus rules.%
\item \href{\url{https://purdue.brightspace.com/d2l/common/assets/pdfjs-d2l-dist/1.0.14-legacy/web/viewer.html?file=\%252Fcontent\%252Fenforced\%252F1095467-wl.202510.CS.24300.001\%252F05L2-ML\%20Basics.pdf\%253Fou\%253D1095467\&lang=en-us\&container=d2l-fileviewer-rendered-pdf\&fullscreen=d2l-fileviewer-rendered-pdf-dialog\&height=667}#page=6}{\textbf{Linear Regression}}: Finding the line that best fits a set of data points, minimizing the sum of squared distances between the points and the line.%
\item \href{\url{https://purdue.brightspace.com/d2l/common/assets/pdfjs-d2l-dist/1.0.14-legacy/web/viewer.html?file=\%252Fcontent\%252Fenforced\%252F1095467-wl.202510.CS.24300.001\%252F05L2-ML\%20Basics.pdf\%253Fou\%253D1095467\&lang=en-us\&container=d2l-fileviewer-rendered-pdf\&fullscreen=d2l-fileviewer-rendered-pdf-dialog\&height=667}#page=11}{\textbf{Multiple Linear Regression}}: Extending linear regression to include multiple independent variables to predict a dependent variable.%
\item \href{\url{https://purdue.brightspace.com/d2l/common/assets/pdfjs-d2l-dist/1.0.14-legacy/web/viewer.html?file=\%252Fcontent\%252Fenforced\%252F1095467-wl.202510.CS.24300.001\%252F05L2-ML\%20Basics.pdf\%253Fou\%253D1095467\&lang=en-us\&container=d2l-fileviewer-rendered-pdf\&fullscreen=d2l-fileviewer-rendered-pdf-dialog\&height=667}#page=13}{\textbf{Curve Fitting}}: Fitting a curve to a set of data points, potentially using polynomial functions of higher degrees.%
\item \href{\url{https://purdue.brightspace.com/d2l/common/assets/pdfjs-d2l-dist/1.0.14-legacy/web/viewer.html?file=\%252Fcontent\%252Fenforced\%252F1095467-wl.202510.CS.24300.001\%252F05L2-ML\%20Basics.pdf\%253Fou\%253D1095467\&lang=en-us\&container=d2l-fileviewer-rendered-pdf\&fullscreen=d2l-fileviewer-rendered-pdf-dialog\&height=667}#page=14}{\textbf{Overfitting}}: The phenomenon where a model fits the training data too closely, resulting in poor generalization to unseen data.%
\item \href{\url{https://purdue.brightspace.com/d2l/common/assets/pdfjs-d2l-dist/1.0.14-legacy/web/viewer.html?file=\%252Fcontent\%252Fenforced\%252F1095467-wl.202510.CS.24300.001\%252F05L2-ML\%20Basics.pdf\%253Fou\%253D1095467\&lang=en-us\&container=d2l-fileviewer-rendered-pdf\&fullscreen=d2l-fileviewer-rendered-pdf-dialog\&height=667}#page=20}{\textbf{Binary Classification}}: Classifying data points into one of two categories using a linear model and a suitable loss function.%
\item \href{\url{https://purdue.brightspace.com/d2l/common/assets/pdfjs-d2l-dist/1.0.14-legacy/web/viewer.html?file=\%252Fcontent\%252Fenforced\%252F1095467-wl.202510.CS.24300.001\%252F05L2-ML\%20Basics.pdf\%253Fou\%253D1095467\&lang=en-us\&container=d2l-fileviewer-rendered-pdf\&fullscreen=d2l-fileviewer-rendered-pdf-dialog\&height=667}#page=25}{\textbf{Choosing a Loss Function}}: Selecting an appropriate loss function (e.g., 0-1 loss, perceptron loss, hinge loss, logistic loss) for binary classification, considering differentiability and computational efficiency.%
\end{itemize}

%
\subsection{Linear Algebra for AI}%
\label{subsec:LinearAlgebraforAI}%
\begin{itemize}[leftmargin=*]%
\item \href{\url{https://purdue.brightspace.com/d2l/common/assets/pdfjs-d2l-dist/1.0.14-legacy/web/viewer.html?file=\%252Fcontent\%252Fenforced\%252F1095467-wl.202510.CS.24300.001\%252F04L1-Linear\%20Algebra\%20Primer.pdf\%253Fou\%253D1095467\&lang=en-us\&container=d2l-fileviewer-rendered-pdf\&fullscreen=d2l-fileviewer-rendered-pdf-dialog\&height=667}#page=5}{\textbf{Representing Linear Transformations with Matrices}}: Matrices are used to represent linear transformations, which map vectors from one space to another while preserving linear combinations. This is fundamental to many AI algorithms.%
\item \href{\url{https://purdue.brightspace.com/d2l/common/assets/pdfjs-d2l-dist/1.0.14-legacy/web/viewer.html?file=\%252Fcontent\%252Fenforced\%252F1095467-wl.202510.CS.24300.001\%252F04L1-Linear\%20Algebra\%20Primer.pdf\%253Fou\%253D1095467\&lang=en-us\&container=d2l-fileviewer-rendered-pdf\&fullscreen=d2l-fileviewer-rendered-pdf-dialog\&height=667}#page=9}{\textbf{Matrix Arithmetic}}: The slide covers basic matrix operations such as addition, subtraction, scalar multiplication, and multiplication, which are essential for many AI computations.%
\item \href{\url{https://purdue.brightspace.com/d2l/common/assets/pdfjs-d2l-dist/1.0.14-legacy/web/viewer.html?file=\%252Fcontent\%252Fenforced\%252F1095467-wl.202510.CS.24300.001\%252F04L1-Linear\%20Algebra\%20Primer.pdf\%253Fou\%253D1095467\&lang=en-us\&container=d2l-fileviewer-rendered-pdf\&fullscreen=d2l-fileviewer-rendered-pdf-dialog\&height=667}#page=11}{\textbf{Solving Systems of Linear Equations using Matrices}}: Systems of linear equations can be represented and solved efficiently using matrix operations, particularly using matrix inverses.%
\item \href{\url{https://purdue.brightspace.com/d2l/common/assets/pdfjs-d2l-dist/1.0.14-legacy/web/viewer.html?file=\%252Fcontent\%252Fenforced\%252F1095467-wl.202510.CS.24300.001\%252F04L1-Linear\%20Algebra\%20Primer.pdf\%253Fou\%253D1095467\&lang=en-us\&container=d2l-fileviewer-rendered-pdf\&fullscreen=d2l-fileviewer-rendered-pdf-dialog\&height=667}#page=16}{\textbf{Matrix Transposition}}: Matrix transposition is a fundamental operation that involves swapping rows and columns of a matrix, useful for data manipulation in AI.%
\item \href{\url{https://purdue.brightspace.com/d2l/common/assets/pdfjs-d2l-dist/1.0.14-legacy/web/viewer.html?file=\%252Fcontent\%252Fenforced\%252F1095467-wl.202510.CS.24300.001\%252F04L1-Linear\%20Algebra\%20Primer.pdf\%253Fou\%253D1095467\&lang=en-us\&container=d2l-fileviewer-rendered-pdf\&fullscreen=d2l-fileviewer-rendered-pdf-dialog\&height=667}#page=18}{\textbf{Parallelizing Matrix Multiplication}}: Domain decomposition is presented as a method for parallelizing matrix multiplication, improving computational efficiency in AI.%
\end{itemize}

%
\section{Algorithm Solutions}%
\label{sec:AlgorithmSolutions}%
\subsection{Linear Algebra}%
\label{subsec:LinearAlgebra}%
\begin{itemize}[leftmargin=*]%
\item \href{\url{https://purdue.brightspace.com/d2l/common/assets/pdfjs-d2l-dist/1.0.14-legacy/web/viewer.html?file=\%252Fcontent\%252Fenforced\%252F1095467-wl.202510.CS.24300.001\%252F04L1-Linear\%20Algebra\%20Primer.pdf\%253Fou\%253D1095467\&lang=en-us\&container=d2l-fileviewer-rendered-pdf\&fullscreen=d2l-fileviewer-rendered-pdf-dialog\&height=667}#page=9}{\textbf{Matrix Addition}}: $A + B = C$, where $c_{ij} = a_{ij} + b_{ij}$%
\item \href{\url{https://purdue.brightspace.com/d2l/common/assets/pdfjs-d2l-dist/1.0.14-legacy/web/viewer.html?file=\%252Fcontent\%252Fenforced\%252F1095467-wl.202510.CS.24300.001\%252F04L1-Linear\%20Algebra\%20Primer.pdf\%253Fou\%253D1095467\&lang=en-us\&container=d2l-fileviewer-rendered-pdf\&fullscreen=d2l-fileviewer-rendered-pdf-dialog\&height=667}#page=9}{\textbf{Matrix Scalar Multiplication}}: $k \cdot A = B$, where $b_{ij} = k \cdot a_{ij}$%
\item \href{\url{https://purdue.brightspace.com/d2l/common/assets/pdfjs-d2l-dist/1.0.14-legacy/web/viewer.html?file=\%252Fcontent\%252Fenforced\%252F1095467-wl.202510.CS.24300.001\%252F04L1-Linear\%20Algebra\%20Primer.pdf\%253Fou\%253D1095467\&lang=en-us\&container=d2l-fileviewer-rendered-pdf\&fullscreen=d2l-fileviewer-rendered-pdf-dialog\&height=667}#page=14}{\textbf{Matrix Multiplication}}: $C_{ij} = \sum_{k=1}^{n} a_{ik}b_{kj}$%
\item \href{\url{https://purdue.brightspace.com/d2l/common/assets/pdfjs-d2l-dist/1.0.14-legacy/web/viewer.html?file=\%252Fcontent\%252Fenforced\%252F1095467-wl.202510.CS.24300.001\%252F04L1-Linear\%20Algebra\%20Primer.pdf\%253Fou\%253D1095467\&lang=en-us\&container=d2l-fileviewer-rendered-pdf\&fullscreen=d2l-fileviewer-rendered-pdf-dialog\&height=667}#page=10}{\textbf{Matrix Inverse}}: $AA^{-1} = A^{-1}A = I$%
\item \href{\url{https://purdue.brightspace.com/d2l/common/assets/pdfjs-d2l-dist/1.0.14-legacy/web/viewer.html?file=\%252Fcontent\%252Fenforced\%252F1095467-wl.202510.CS.24300.001\%252F04L1-Linear\%20Algebra\%20Primer.pdf\%253Fou\%253D1095467\&lang=en-us\&container=d2l-fileviewer-rendered-pdf\&fullscreen=d2l-fileviewer-rendered-pdf-dialog\&height=667}#page=12}{\textbf{Solving Linear Equations using Matrix Inverses}}: $Ax = b \implies x = A^{-1}b$%
\end{itemize}

%
\subsection{Calculus}%
\label{subsec:Calculus}%
\begin{itemize}[leftmargin=*]%
\item \href{\url{https://purdue.brightspace.com/d2l/common/assets/pdfjs-d2l-dist/1.0.14-legacy/web/viewer.html?file=\%252Fcontent\%252Fenforced\%252F1095467-wl.202510.CS.24300.001\%252F04L1-Linear\%20Algebra\%20Primer.pdf\%253Fou\%253D1095467\&lang=en-us\&container=d2l-fileviewer-rendered-pdf\&fullscreen=d2l-fileviewer-rendered-pdf-dialog\&height=667}#page=21}{\textbf{Derivative of Sigmoid Function}}: $\frac{d}{dx} \sigma(x) = \frac{e^{-x}}{(1 + e^{-x})^2}$%
\end{itemize}

%
\subsection{Machine Learning Algorithms}%
\label{subsec:MachineLearningAlgorithms}%
\begin{itemize}[leftmargin=*]%
\item \href{\url{https://purdue.brightspace.com/d2l/common/assets/pdfjs-d2l-dist/1.0.14-legacy/web/viewer.html?file=\%252Fcontent\%252Fenforced\%252F1095467-wl.202510.CS.24300.001\%252F05L2-ML\%20Basics.pdf\%253Fou\%253D1095467\&lang=en-us\&container=d2l-fileviewer-rendered-pdf\&fullscreen=d2l-fileviewer-rendered-pdf-dialog\&height=667}#page=30}{\textbf{Perceptron Algorithm}}: Iteratively update the weight vector $w$ based on misclassified data points. If $sign(w^Tx_n) \ne y_n$, then $w \leftarrow w + y_nx_n$.%
\end{itemize}

%
\subsection{Search Algorithms}%
\label{subsec:SearchAlgorithms}%
\begin{itemize}[leftmargin=*]%
\item \href{\url{https://purdue.brightspace.com/d2l/common/assets/pdfjs-d2l-dist/1.0.14-legacy/web/viewer.html?file=\%252Fcontent\%252Fenforced\%252F1095467-wl.202510.CS.24300.001\%252F02L2-Intro.pdf\%253Fou\%253D1095467\&lang=en-us\&container=d2l-fileviewer-rendered-pdf\&fullscreen=d2l-fileviewer-rendered-pdf-dialog\&height=667}#page=22}{\textbf{Minimax Search}}: A search algorithm for two-player zero-sum games that aims to find the best move for the current player, assuming the opponent also plays optimally. It uses an evaluation function to assign scores to game states and recursively explores the game tree, alternating between maximizing the score for the current player and minimizing the score for the opponent.%
\item \href{\url{https://purdue.brightspace.com/d2l/common/assets/pdfjs-d2l-dist/1.0.14-legacy/web/viewer.html?file=\%252Fcontent\%252Fenforced\%252F1095467-wl.202510.CS.24300.001\%252F02L2-Intro.pdf\%253Fou\%253D1095467\&lang=en-us\&container=d2l-fileviewer-rendered-pdf\&fullscreen=d2l-fileviewer-rendered-pdf-dialog\&height=667}#page=23}{\textbf{Alpha-Beta Pruning}}: A technique to reduce the search space in Minimax search by avoiding the evaluation of branches that cannot affect the optimal decision. It uses alpha (best value for the maximizing player) and beta (best value for the minimizing player) values to prune branches.%
\end{itemize}

%
\end{document}